\documentclass[12pt,a4paper]{article}
\usepackage{fontspec}
\usepackage{polyglossia}
\setdefaultlanguage{russian}
\setmainfont{Arial}
\setmonofont{Arial}
\newfontfamily\cyrillicfonttt{Arial}
\usepackage[left=2cm,right=2cm,top=2cm,bottom=2cm]{geometry}
\usepackage{graphicx,booktabs,float,hyperref,xurl,fancyhdr,amsmath}
\hypersetup{colorlinks=true,urlcolor=blue}
\pagestyle{fancy}\fancyhf{}
\lhead{Поиск ассоциативных правил --- Apriori}\rhead{\thepage}
\setlength{\headheight}{15pt}

\begin{document}
\begin{titlepage}\centering
\vspace*{3cm}{\Huge\bfseries ОТЧЁТ О ВЫПОЛНЕНИИ ЗАДАНИЯ\par}
\vspace{1cm}{\LARGE\bfseries Поиск ассоциативных правил\\на основе алгоритма Apriori\par}
\vspace{3cm}\fbox{\begin{minipage}{0.78\textwidth}
Набор данных: \textbf{baskets.csv}\\
Число транзакций: \textbf{7501}\\
Минимальная поддержка: \textbf{1\%}\\
Пороги достоверности: \textbf{25\%, 30\%, 35\%, 40\%, 45\%, 50\%}
\end{minipage}}\vfill 2026
\end{titlepage}

\section{Формулировка задания}
Необходимо дополнить программу поиска частых наборов построением ассоциативных правил в форме «антецедент $\to$ консеквент», рассчитывать поддержку и достоверность, поддержать порог достоверности и сортировку по поддержке либо лексикографически. Требуется исследовать время поиска и число правил, построить диаграммы и интерпретировать правила суммарной длиной не более семи объектов.

\section{Материалы проекта}
Исходный код, данные и результаты: \url{https://github.com/FeodorG/data-mining-fundamentals/tree/main/lab2}.

\section{Методика}
После поиска частых наборов Apriori для каждого набора $X$ длиной не менее двух перебираются непустые собственные подмножества $A$, образующие правило $A\to X\setminus A$. Поддержка равна доле транзакций, содержащих обе части правила, а достоверность вычисляется как $\operatorname{support}(A,B)/\operatorname{support}(A)$.

При поддержке 1\% максимальная достоверность равна 50{,}67\%. Поэтому диапазон 70--95\% дал бы только нулевые результаты; для содержательного эксперимента выбран диапазон 25--50\% с шагом 5\%. Время генерации правил измерялось 100 раз, далее приведена медиана.

\section{Результаты экспериментов}
\begin{table}[H]\centering
\caption{Результаты при фиксированной поддержке 1\%}
\begin{tabular}{rrr}\toprule
Достоверность & Медиана, мс & Правил\\\midrule
25\% & 0{,}228 & 99\\30\% & 0{,}214 & 65\\35\% & 0{,}200 & 34\\
40\% & 0{,}196 & 19\\45\% & 0{,}190 & 6\\50\% & 0{,}187 & 2\\\bottomrule
\end{tabular}\end{table}

\begin{figure}[H]\centering
\includegraphics[width=.93\textwidth]{results/rules_runtime.png}
\caption{Время генерации и фильтрации правил}
\end{figure}
\begin{figure}[H]\centering
\includegraphics[width=.93\textwidth]{results/rules_count.png}
\caption{Число правил при изменении порога достоверности}
\end{figure}

Число правил монотонно сокращается с 99 до 2. Время уменьшается лишь немного, поскольку набор перебираемых кандидатов остаётся прежним, а меняется главным образом число добавляемых результатов.

\section{Правила длиной не более семи объектов}
Максимальная длина частого набора равна трём, поэтому все 99 правил при достоверности 25\% удовлетворяют ограничению. Полный список сохранён в \path{results/rules_upto7.csv}. Самые достоверные примеры:
\begin{itemize}
\item \{говяжий фарш, яйца\} $\to$ \{минеральная вода\}: поддержка 1{,}01\%, достоверность 50{,}67\%;
\item \{говяжий фарш, молоко\} $\to$ \{минеральная вода\}: 1{,}11\% и 50{,}30\%;
\item \{суп\} $\to$ \{минеральная вода\}: 2{,}31\% и 45{,}65\%.
\end{itemize}

Минеральная вода часто выступает консеквентом и может рассматриваться как кандидат для совместной выкладки или рекомендации. Направление правила важно: для «макароны $\to$ минеральная вода» достоверность около 32{,}6\%, а для обратного --- около 25{,}7\% при одинаковой поддержке. Правила показывают совместную встречаемость, но не доказывают причинную связь.

\section{Вывод}
Программа дополнена генерацией правил, читаемым выводом, CSV-экспортом, двумя видами сортировки и ограничением размера. Эксперимент показал, что повышение достоверности резко сокращает число правил, но мало влияет на время их генерации.
\end{document}
